% CV de Roddy Edison Huarhua Rojas
% Compilar con: pdflatex cv_roddy_huarhua.tex
\documentclass[10pt,a4paper]{article}
\usepackage[utf8]{inputenc}
\usepackage[T1]{fontenc}
\usepackage[spanish]{babel}
\usepackage{lmodern}
\usepackage[sfdefault]{roboto}
\usepackage[a4paper,margin=1.7cm,top=1.5cm,bottom=1.5cm]{geometry}
\usepackage{xcolor}
\usepackage{titlesec}
\usepackage{enumitem}
\usepackage{fontawesome5}
\usepackage[hidelinks]{hyperref}
\usepackage{tabularx}
\usepackage{fancyhdr}
\usepackage{needspace}

\definecolor{teal}{HTML}{1F5F5B}
\definecolor{terra}{HTML}{A8492A}
\definecolor{muted}{HTML}{5B6669}

\hypersetup{pdftitle={CV - Roddy Huarhua Rojas},pdfauthor={Roddy Edison Huarhua Rojas}}
\setlength{\parindent}{0pt}
\setlength{\parskip}{0pt}
\pagestyle{fancy}\fancyhf{}\renewcommand{\headrulewidth}{0pt}
\fancyfoot[R]{\small\color{muted}Roddy E. Huarhua Rojas · \thepage}

\titleformat{\section}{\large\bfseries\color{teal}}{}{0pt}{\MakeUppercase}[\vspace{-6pt}{\color{terra}\rule{\linewidth}{0.8pt}}]
\titlespacing*{\section}{0pt}{12pt}{6pt}
\setlist[itemize]{leftmargin=1.2em,itemsep=1.5pt,topsep=2pt,label={\color{terra}\small$\bullet$}}

% \trabajo{Empresa}{Cargo}{Fechas}{Lugar}
\newcommand{\trabajo}[4]{%
  \vspace{5pt}\textbf{#1} \hfill {\small\color{muted}#3}\\
  {\color{teal}\itshape #2}\hfill{\small\color{muted}#4}\par\vspace{1pt}}

\begin{document}

% ---------- Encabezado ----------
{\fontsize{24}{28}\selectfont\bfseries Roddy Edison Huarhua Rojas}\\[3pt]
{\large\color{teal} Gestión estratégica de relaciones comunitarias · Minería, energía e infraestructura}\\[6pt]
{\small\color{muted}
\faMapMarker*\ Lima, Perú \quad
\faEnvelope\ \href{mailto:huarhuarojas@gmail.com}{huarhuarojas@gmail.com} \quad
\faLinkedin\ \href{https://www.linkedin.com/in/roddy-huarhua/}{linkedin.com/in/roddy-huarhua}\\[2pt]
\faGlobe\ \href{https://roddyhuarhua-hub.github.io}{roddyhuarhua-hub.github.io} \quad
\faGithub\ \href{https://github.com/roddyhuarhua-hub}{roddyhuarhua-hub}}

\section{Perfil ejecutivo}
Profesional de ciencias sociales, licenciado en Historia, con más de 14 años de trayectoria en relaciones comunitarias, gestión social y prevención y transformación de conflictos socioambientales en minería, energía, gas natural e infraestructura. Experiencia coordinando estrategias de relacionamiento territorial, negociaciones de acceso a tierras y servidumbres, participación ciudadana, cumplimiento de compromisos y programas de desarrollo social. Capacidad para articular comunidades campesinas, autoridades, organismos públicos, equipos operativos, consultoras y contratistas, con orientación a la continuidad de los proyectos y a la construcción de acuerdos sostenibles.

\section{Resultados relevantes}
\begin{itemize}
  \item Gestión social de \textbf{ocho proyectos} de exploración minera y obtención de \textbf{seis permisos sociales} durante 2025 en Buenaventura.
  \item Participación en la gestión de \textbf{tres paralizaciones} de obra, negociación con \textbf{35 posesionarios} y cierre de más de \textbf{30 observaciones sociales} en un proyecto de transmisión eléctrica de 220 kV.
  \item Negociación de una \textbf{servidumbre comunal por seis años} para dar continuidad a cuatro campañas de perforación diamantina en Palcawanka.
  \item Diseño y ejecución de \textbf{Cuentigas}, iniciativa educativa desarrollada durante tres años en más de 50 instituciones educativas de Ica, con más de 200 relatos estudiantiles presentados.
\end{itemize}

\section{Competencias directivas y técnicas}
Estrategia de relaciones comunitarias · Diálogo intercultural y negociación · Gestión de conflictos socioambientales · Participación ciudadana (DIA/EIA) · Servidumbres y acuerdos prediales · Gestión de compromisos y reclamos · Inversión social y desarrollo local · Mapeo y análisis de actores · Gestión de contratistas y proveedores locales · Articulación con autoridades y entidades públicas · Comunicación externa y relacionamiento con prensa · Reportes ejecutivos e indicadores sociales · \textbf{Análisis de datos con Python} (pandas, mapas, dashboards)

\section{Experiencia profesional}
\trabajo{Compañía de Minas Buenaventura S.A.A.}{Coordinador Local de Asuntos Sociales}{nov. 2024 – actualidad}{Pasco}
\begin{itemize}
  \item Coordino la gestión social de ocho proyectos de exploración, la negociación de permisos y el relacionamiento con comunidades, autoridades y organizaciones del área de influencia.
  \item Conduzco espacios de diálogo y asambleas, mapeo de actores, seguimiento de compromisos, identificación de riesgos y coordinación con las áreas de Exploraciones, Legal, Ambiental y Operaciones.
  \item Resultados: seis permisos sociales obtenidos en 2025; más de 50 oportunidades de empleo directo e indirecto vinculadas al proyecto Anaís; articulación de servicios con proveedores y empresas comunales.
  \item Participación en negociación de servidumbre de seis años para el proyecto Sami, en proceso de formalización; coordinación de iniciativas de infraestructura y fondos sociales comunitarios.
\end{itemize}

\trabajo{Conelsur}{Gestor Social · Proyecto de transmisión 220 kV}{ago. 2023 – nov. 2024}{}
\begin{itemize}
  \item Gestioné el relacionamiento comunitario durante la construcción de una línea de transmisión eléctrica, incluyendo posesionarios, comunidades campesinas, autoridades y contratistas.
  \item Participé en mesas de diálogo con intervención del MINEM y en procesos de negociación para superar controversias y mantener la continuidad de las obras.
  \item Resultados: contribución al levantamiento de tres paralizaciones, acuerdos con 35 posesionarios y cierre de más de 30 observaciones sociales del proceso constructivo.
\end{itemize}

\trabajo{Consorcio Minero Palcawanka S.A.C.}{Coordinador de Relaciones Comunitarias}{nov. 2020 – ago. 2023}{Huancavelica}
\begin{itemize}
  \item Coordiné la estrategia social para exploración minera y el diálogo con comunidades campesinas, dirigentes y autoridades; prevención de conflictos y seguimiento de acuerdos.
  \item Negocié una servidumbre comunal por seis años que permitió la continuidad de cuatro campañas de perforación diamantina; promoví empleo rotativo, compras y servicios locales.
\end{itemize}

\trabajo{LQP Predial y Social S.A.C.}{Especialista Social Senior}{feb. – oct. 2020}{Ayacucho}
\begin{itemize}
  \item Desarrollé mapeo de actores, estrategias de ingreso social, reuniones informativas y articulación con anexos y autoridades locales para actividades de exploración.
\end{itemize}

\trabajo{Río Blanco Copper S.A.}{Coordinador de Relaciones Comunitarias}{jul. 2018 – ene. 2020}{Piura}
\begin{itemize}
  \item Coordiné relacionamiento en un entorno de alta sensibilidad social; desarrollé redes de aliados locales, seguimiento de riesgos, diálogo con organizaciones territoriales y apoyo a iniciativas educativas y productivas.
  \item Contribuí al ingreso y relacionamiento en el área de influencia indirecta del proyecto mediante acercamiento con actores y autoridades.
\end{itemize}

\trabajo{Contugas S.A.C.}{Analista Senior de Relaciones Comunitarias}{abr. 2015 – jul. 2018}{Ica}
\begin{itemize}
  \item Implementé acciones de relaciones comunitarias y compromisos sociales asociados a instrumentos ambientales en Chincha, Pisco, Ica, Nasca y Marcona; gestioné consultas, reclamos y coordinación con contratistas.
  \item Diseñé e impulsé programas de educación y comunicación social, incluido Cuentigas (más de 50 colegios y 200 relatos), en coordinación con UGEL y Dirección Regional de Educación de Ica.
  \item Participé en comunicación externa y relacionamiento con medios de prensa, así como en el acompañamiento de producciones audiovisuales.
\end{itemize}

\trabajo{Consorcio Minero Horizonte}{Relaciones Comunitarias / Consultor}{ene. 2012 – jun. 2013 y jun. 2014 – mar. 2015}{Huancavelica}
\begin{itemize}
  \item Negociación de servidumbres y compensaciones, participación ciudadana en exploración y prevención de conflictos; implementación de empleo local rotativo.
  \item En la primera etapa, coordinación operativa de 24 trabajadores manuales y cuatro técnicos, con control de asistencia, logística y rendiciones.
\end{itemize}

\trabajo{Ministerio de Vivienda}{Gestor Institucional, Programa Nacional Tambos}{may. – jul. 2014}{Huancavelica}
\begin{itemize}
  \item Articulación entre instituciones públicas, gobiernos locales y población para canalizar servicios y necesidades territoriales.
\end{itemize}

\trabajo{Desafío Minero S.A.C.}{Consultor de Relaciones Comunitarias}{may. – dic. 2011}{Huancavelica y Arequipa}
\begin{itemize}
  \item Apoyo a negociación, talleres de participación ciudadana, prevención de conflictos y relacionamiento con comunidades.
\end{itemize}

\section{Liderazgo empresarial y consultoría}
\trabajo{Confluentia S.A.C.}{Gerente General y Fundador}{ago. 2026 – actualidad}{\href{https://www.confluentia.pe}{confluentia.pe}}
Dirección de iniciativa de consultoría en gestión social, negociación, prevención de conflictos y acompañamiento territorial. Actividad paralela a la posición actual.

\trabajo{Mekary Consulting S.A.C.}{Principal CEO}{jul. 2022 – ene. 2025}{}
Dirección estratégica y desarrollo comercial de servicios de consultoría social y ambiental; coordinación de propuestas, equipos y proyectos.

\needspace{8\baselineskip}
\section{Formación académica}
\trabajo{Pontificia Universidad Católica del Perú}{Maestría en Gerencia Social, mención en Gerencia del Desarrollo Económico Local}{2020 – 2022}{}
Tesis: \href{https://hdl.handle.net/20.500.12404/28120}{\itshape Percepciones y estrategias de participación en entornos de variación socioeconómica. Caso del proyecto de exploración Palcawanka, comunidad campesina de Conaicasa} (2024).

\trabajo{Universidad Nacional Federico Villarreal}{Licenciado en Historia}{}{}
Tesis: \href{https://hdl.handle.net/20.500.13084/6906}{\itshape Comunidad campesina de Palca y la empresa Consorcio Minero Horizonte: entre el conflicto socioambiental y el diálogo (2010–2011)} (2023).

\trabajo{Pontificia Universidad Católica del Perú}{Diplomatura de Especialización en Ciencia de Datos para las Ciencias Sociales y Gestión Pública}{ago. 2026 – may. 2027 (en curso)}{}
Portafolio de proyectos con Python: \href{https://roddyhuarhua-hub.github.io}{roddyhuarhua-hub.github.io}

\section{Especializaciones}
\begin{tabularx}{\linewidth}{@{}X r@{}}
Transformación de Conflictos Socioambientales para el Desarrollo — Universidad Antonio Ruiz de Montoya & {\color{muted}2015}\\
Diseño y Gestión de Proyectos Sociales — PUCP & {\color{muted}2018–2019}\\
Gerencia Social — PUCP & {\color{muted}2020–2021}\\
Sostenibilidad Social y Comunicación Estratégica — Universidad Científica del Sur & {\color{muted}2023–2024}\\
Administración — Zegel & {\color{muted}2023}\\
\end{tabularx}

\section{Información adicional}
\textbf{Idiomas:} español nativo; inglés en formación intermedia.\\
\textbf{Herramientas:} Microsoft Office; Python (pandas, geopandas, Streamlit) y R, en formación.\\
\textbf{Otros:} licencia de conducir A-I vigente. Disponibilidad para trabajo de campo y desplazamientos.

\end{document}
